\documentclass[letterpaper,11pt]{moderncv}
\moderncvstyle{banking}
\moderncvcolor{black}

\usepackage[letterpaper, margin=0.8in]{geometry}
\usepackage{xcolor}
\usepackage{fontspec}
\setmainfont{Times New Roman}
\usepackage[unicode]{hyperref}

\name{Saeyoung}{Kim}
\address{Cambridge, MA, USA}{}{}
\phone[mobile]{+1 (351) 322 5510}
\email{saeyoung\_kim@uml.edu}
\social[linkedin]{saeyoung-macx-kim}
\social[github]{PrimaryAnomaly}

\recipient{Hiring Team}{Figure\\San Jose, CA}
\date{\today}
\opening{Dear Hiring Team,}
\closing{Sincerely,}

\begin{document}

\makelettertitle

I am applying for the Reliability Test Engineer, Hardware position at Figure. I have 9+ years of experience designing reliability test setups, writing test plans, and running root cause analysis on electromechanical systems from prototype through production.

At Hanwha Systems, I designed reliability test setups for satellite subsystems based on mission life requirements: power distribution, processors, sensors, and fully integrated systems. I wrote and authorized test procedures, trained technicians, and ran multiple parallel test campaigns. When hardware failed during integration, I led the root cause investigations, developed validation protocols to confirm fixes, and reported results across engineering teams. I designed test fixtures in CAD with GD\&T and worked with suppliers on custom test equipment.

During my PhD, I validated multi-modal sensor suites across thermal cycling and environmental stress testing. I built DAQ systems with accelerometers, thermal probes, and electrochemical sensors, and took hardware prototypes from design through production-level testing. In my postdoc, I run reliability test plans for manufacturing systems, oversee parallel test initiatives, and lead root cause analysis when things go wrong.

A humanoid robot that works in the real world has to be reliable, and proving that reliability requires serious test engineering. That is what I do, and I would like to do it at Figure.

Thank you for your consideration.

\makeletterclosing

\end{document}
